% =============================================================================
% METADATOS DEL DOCUMENTO (plantilla genérica para prácticas de laboratorio)
% Este archivo debe cargarse AL PRINCIPIO del preámbulo.
% =============================================================================
%
% INSTRUCCIONES: Reemplaza los textos entre corchetes por los valores reales
% de tu compendio. Puedes eliminar los comentarios que no necesites.
% =============================================================================

% --- Identificación del documento ---
\newcommand{\docTitulo}{[Título del compendio de prácticas]}           % Ej: "Prácticas de Laboratorio de Física Computacional"
\newcommand{\docSubtitulo}{[Subtítulo o descripción breve]}           % Opcional: puede dejarse vacío
\newcommand{\docAutor}{[Nombre del autor o autores]}                  % Ej: "J. A. Chala-Casanova"
\newcommand{\docFecha}{\the\year}                                     % Año actual (o fecha fija)
\newcommand{\docEstado}{[Versión]}                                    % Ej: "Borrador", "Versión 1.0"

% --- Cita de portada (opcional) ---
\newcommand{\docCita}{%
  ``[Cita inspiradora relacionada con el tema del compendio]''%
}

% --- Textos para páginas preliminares (opcionales) ---
\newcommand{\docDedicatoria}{}      % Si no se usa, dejar vacío
\newcommand{\docEpigrafe}{}         % Si no se usa, dejar vacío
\newcommand{\docAgradecimientos}{}  % Si no se usa, dejar vacío

% --- Palabras clave y clasificación (para metadatos PDF) ---
\newcommand{\docKeywords}{[palabra clave 1, palabra clave 2, ...]}   % Ej: "física computacional, machine learning"
\newcommand{\docSubject}{[Área temática principal]}                   % Ej: "Física Computacional / Machine Learning"

% =============================================================================
% NUEVOS COMANDOS (opcionales, pueden usarse en la portada o en secciones)
% =============================================================================

% --- Afiliación del autor ---
\newcommand{\docAfiliacion}{[Institución, departamento]}              % Ej: "Departamento de Física, Universidad X"

% --- Fechas específicas (si se desea distinguir realización y entrega) ---
\newcommand{\docFechaRealizacion}{\today}
\newcommand{\docFechaEntrega}{\today}

% --- Resumen general del compendio ---
\newcommand{\docAbstract}{%
  [Breve descripción del contenido del compendio, público objetivo y metodología.]
}

% --- Referencia a material complementario (opcional) ---
\newcommand{\docLibroReferencia}{%
  [Información sobre material asociado, como un libro o curso.]
}


% =============================================================================
% RECURSOS BIBLIOGRÁFICOS
% =============================================================================
\addbibresource{references/references.bib}

